\section{A finite criterion on directed edges}\label{dt:criterion}
The root in this section is an order for evaluating the closed finite
contraction. It need not be the terminal root used to draw the original
knot. A nonroot vertex sends one polynomial to its parent; the last
vertex sends none. The following constants and inequalities reflect
that difference.

Put

\begin{equation}\label{dt:eq-1}
\begin{gathered}
a_0=\frac1{2\pi},\qquad k_0=2^{-16},\qquad
 R_0=\frac18,\qquad b_0=2^{-11},\qquad\\
 r_0=\frac1{16},\qquad t_0=2^{-14},\qquad
 \chi=1-\cos r_0.
\end{gathered}
\end{equation}

The parent-variable rectangle is

\begin{equation}\label{dt:eq-2}
\mathcal R=\{u:|\Re u|\le R_0,\ |\Im u|\le b_0\}.
\end{equation}

Choose any vertex as root and direct all edges toward it. For every
nonroot vertex choose a number

\begin{equation}\label{dt:eq-3}
0<c_v<k_0,
 \qquad
 \delta_v=\min\left\{\frac{b_0}2,
               \frac{c_v\chi}{4(1+c_v)}\right\}.
\end{equation}

These are parameters of the criterion. They can differ at every edge.
They are not asymptotic quantities to be inferred from the invariant.

Write \(q_v=1\) at a nonroot vertex and \(q_v=0\) at the root. If \(\mathrm{Ch}(v)\) is the
set of children, define

\begin{equation}\label{dt:eq-4}
\begin{aligned}
 a_v&=\frac{d_v}{2\pi},& b_v&=\frac{n_v}{2\pi},&
 D_v&=a_v-i b_v,\\
 A_v&=\frac2{\pi^2}
       \left(q_vk_0+\frac12\sum_{w\in\mathrm{Ch}(v)}c_w\right),&
 B_v&=16q_v+|\mathrm{Ch}(v)|,&
 C_v&=B_v+\frac{b_v^2}{4A_v}.
\end{aligned}
\end{equation}

There is at least one incidence at every vertex, so \(A_v,B_v\) are positive.
Let

\begin{equation}\label{dt:eq-5}
\delta_v^{\mathrm{in}}
      =\min\bigl(\{b_0\}\cup
                     \{\delta_w:w\in\mathrm{Ch}(v)\}\bigr)
\end{equation}

and set

\begin{equation}\label{dt:eq-6}
y_v=\min\left\{
 \frac{\delta_v^{\mathrm{in}}}2,
 \frac{t_0|D_v|}{4C_v},
 \frac{A_v\pi^2}{8(\pi|b_v|+t_0|D_v|)},
 \sqrt{\frac{A_v\pi^2}{8B_v}}
 \right\}.
\end{equation}

Finally define two action quantities

\begin{equation}\label{dt:eq-7}
\beta_v=\frac{210}{361}\frac{a_v}{|D_v|^2},
 \qquad
 \eta_v=\min\left\{
 \frac12t_0|D_v|y_v,
 \frac12A_v\pi^2,
 A_vr_0^2,
 \frac18a_vr_0^2
 \right\}.
\end{equation}

The criterion consists of the strict inequalities

\begin{equation}\label{dt:eq-8}
\boxed{\quad t_0|D_v|\ge2,\qquad
                 \beta_v<\frac{\eta_v}{2}\quad
                      \text{at every vertex},\quad}
\end{equation}

and, only at nonroot vertices,

\begin{equation}\label{dt:eq-9}
\boxed{\qquad
       \beta_v<(k_0-c_v)\chi.\qquad}
\end{equation}

It is a finite elementary test in the degrees, signed powers, root and
positive numbers \(c_v\). It uses no value of an unknown quantum invariant
and no assumed dominance of a global saddle. Its quantities are local:
\(A_v\) and the available strip depend on the child budgets, while the exact
child scalar will be divided out rather than estimated.


We say that the chosen root and numbers \(c_v\) satisfy the directed
criterion when \eqref{dt:eq-3}, \eqref{dt:eq-8} and \eqref{dt:eq-9} hold.
No numerical estimate of an invariant enters this definition.
The number \(\beta_v\) bounds the action lost in evaluating just the
current trace after exact child constants are divided out. The four
entries of \(\eta_v\) pay, respectively, for the infinite charge tail,
the original endpoints, localization to a real interval, and its
vertical connecting paths. The extra comparison \eqref{dt:eq-9} is used
only to produce an outgoing real bound.
